% ================================================================== Chapitre 4
\chapter{Premier principe : effets thermiques des réactions}
\label{ch:premier}

Une réaction chimique libère ou absorbe de l'énergie. Selon que le réacteur échange de la chaleur avec l'extérieur ou non, cette énergie se manifeste par un \emph{transfert thermique} ou par une \emph{variation de température}. L'outil est le premier principe, écrit avec l'enthalpie.

\section{Vocabulaire des transformations}

\begin{definition}[conditions opératoires]
Pour une transformation entre un état initial (I) et un état final (F) :
\begin{itemize}
  \item \textbf{monobare} : la pression extérieure est constante et $P_I = P_F = P_{\text{ext}}$ (réacteur ouvert sur l'atmosphère) ;
  \item \textbf{isobare} : la pression du système est définie et constante à chaque instant ;
  \item \textbf{monotherme} : le système n'échange de chaleur qu'avec un thermostat de température $T_{\text{ext}}$ ;
  \item \textbf{isotherme} : $T$ est définie et constante tout au long de la transformation (en pratique : $T_I = T_F = T_{\text{ext}}$ pour un réacteur thermostaté) ;
  \item \textbf{adiabatique} : aucun transfert thermique ($Q = 0$), par exemple dans un calorimètre ou pour une réaction très rapide.
\end{itemize}
\end{definition}

\begin{rappel}[premier principe en enthalpie]
Pour un système fermé, au repos macroscopique, subissant une transformation \emph{monobare} où seules les forces de pression travaillent :
\[ \Delta H = Q \]
($Q$ : transfert thermique \emph{reçu} par le système, algébrique).
\end{rappel}

\section{Variation d'enthalpie due à la réaction}

À $T$ et $P$ fixées, $\dd H = \Delta_rH\,\dd\xi$. On utilise ensuite l'approximation suivante, toujours faite au programme :
\[ \Delta_rH \simeq \Delta_rH^\circ. \]
Elle est justifiée car l'enthalpie d'un gaz parfait ne dépend que de $T$ (pas de sa pression partielle), celle d'une phase condensée dépend très peu de $P$, et le mélange idéal se fait sans effet thermique. Ainsi, pour une réaction menée de $\xi = 0$ à $\xi_F$ à température constante :
\[ \boxed{\Delta H = \xi_F\,\Delta_rH^\circ(T)} \]

\section{Réacteur monobare et isotherme : transfert thermique}

\begin{resultat}{Transfert thermique en réacteur monobare isotherme}
\[ \boxed{Q = \Delta H = \xi_F\,\Delta_rH^\circ} \]
\begin{itemize}
  \item réaction \textbf{exothermique} ($\Delta_rH^\circ<0$) : $Q<0$, le système \emph{cède} de la chaleur au milieu extérieur ;
  \item réaction \textbf{endothermique} ($\Delta_rH^\circ>0$) : $Q>0$, le système \emph{reçoit} de la chaleur.
\end{itemize}
\end{resultat}

\begin{exemple}[chauffage au gaz naturel]
On brûle $\SI{1}{m^3}$ de méthane (pris à $\SI{1}{bar}$ et $\SI{298}{K}$) : $n = \dfrac{PV}{RT} = \dfrac{10^5\times1}{8{,}314\times298} = \SI{40,3}{mol}$.

\noindent Réaction : \ce{CH4(g) + 2 O2(g) -> CO2(g) + 2 H2O(l)}, avec (loi de Hess)
\[ \Delta_rH^\circ = -393{,}5 + 2(-285{,}8) - (-74{,}6) = \SI{-890,5}{kJ.mol^{-1}}. \]
Si les produits sont ramenés à $\SI{298}{K}$ (eau condensée) : $Q = \xi_F\Delta_rH^\circ = 40{,}3\times(-890{,}5) = \SI{-35,9}{MJ} \approx \SI{-10}{kWh}$. C'est l'ordre de grandeur du « pouvoir calorifique » facturé par les fournisseurs de gaz ($\approx\SI{10}{kWh}$ par $\si{m^3}$).
\end{exemple}

\section{Réacteur monobare et adiabatique : température finale}

Si le réacteur est calorifugé (ou la réaction très rapide), $Q = 0$, donc $\Delta H = 0$ : l'énergie libérée par une réaction exothermique \emph{chauffe} le contenu du réacteur.

Pour calculer $T_F$, on remplace la transformation réelle (réaction et échauffement simultanés) par un \textbf{chemin fictif} en deux étapes, ce qui est légitime puisque $H$ est une fonction d'état.

\begin{figure}[H]
\centering
\begin{tikzpicture}[>=Stealth,box/.style={draw=navy,line width=0.8pt,rounded corners=2pt,fill=cdefbg,inner sep=5pt,align=center,font=\small,minimum width=3.6cm}]
  \node[box] (I) at (0,0) {État initial\\réactifs, $T_I$, $P^\circ$};
  \node[box] (A) at (8,0) {État intermédiaire\\système final, $T_I$, $P^\circ$};
  \node[box,fill=cresbg,draw=c3] (F) at (8,3.4) {État final\\système final, $T_F$, $P^\circ$};
  \draw[->,line width=1.2pt,c3,dashed] (I.north) -- node[above left,font=\small,align=center] {transformation réelle\\adiabatique : $\Delta H = 0$} (F.west);
  \draw[->,line width=1.1pt,navy] (I) -- node[above,font=\small] {(a) réaction à $T_I$} node[below,font=\small] {$\Delta H_a = \xi_F\,\Delta_rH^\circ(T_I)$} (A);
  \draw[->,line width=1.1pt,navy] (A) -- node[right,font=\small,align=left] {(b) échauffement :\\$\Delta H_b = \Big(\sum_j n_jC^\circ_{p,m,j}\Big)(T_F-T_I)$} (F);
\end{tikzpicture}
\caption{Chemin fictif pour le calcul d'une température finale en réacteur adiabatique. La somme porte sur \emph{tous} les constituants présents dans l'état final.}
\label{fig:cheminfictif}
\end{figure}

\begin{resultat}{Température finale en réacteur adiabatique monobare}
$\Delta H_a + \Delta H_b = 0$, d'où, sans changement d'état au cours de l'échauffement :
\[ \boxed{T_F = T_I - \frac{\xi_F\,\Delta_rH^\circ(T_I)}{\sum_j n_j\,C^\circ_{p,m,j}}} \]
\end{resultat}

\begin{methode}[température de fin de réaction en réacteur adiabatique]
\begin{enumerate}
  \item Écrire la réaction et dresser le tableau d'avancement ; trouver $\xi_F$ (souvent $\xi_{\max}$ pour une réaction totale).
  \item Faire l'inventaire de l'état final : produits, réactifs en excès \emph{et espèces spectatrices} (le diazote de l'air, le solvant\dots).
  \item Écrire $\Delta H = \Delta H_a + \Delta H_b = 0$ le long du chemin fictif.
  \item Si un constituant change d'état pendant l'échauffement, couper l'étape (b) à la température de changement d'état et ajouter $n\,\Delta_{\text{vap}}H^\circ$ (ou $\Delta_{\text{fus}}H^\circ$).
  \item Vérifier la cohérence : $T_F > T_I$ pour une réaction exothermique.
\end{enumerate}
\end{methode}

\begin{definition}[température de flamme adiabatique]
La \textbf{température de flamme} d'une combustion est la température finale atteinte par les gaz brûlés lorsque la combustion, totale, a lieu dans des conditions adiabatiques et monobares.
\end{definition}

\begin{exemple}[flamme méthane--air]
Mélange stœchiométrique dans l'air (environ 4 \ce{N2} pour 1 \ce{O2}) à $T_I = \SI{298}{K}$ : on part de \ce{CH4 + 2 O2 + 8 N2}, et l'état final est \ce{CO2 + 2 H2O(g) + 8 N2}. On choisit dans l'étape (a) de produire directement de l'eau \emph{gazeuse} (état standard gaz à $\SI{298}{K}$, fictif mais légitime), ce qui évite de traiter sa vaporisation :
\begin{gather*} \Delta_rH^\circ = -393{,}5 + 2(-241{,}8) + 74{,}6 = \SI{-802,5}{kJ.mol^{-1}},\\ \sum_jn_jC^\circ_{p,m,j} = 37{,}1 + 2(33{,}6) + 8(29{,}1) = \SI{337,1}{J.K^{-1}}. \end{gather*}
\[ T_F = 298 + \frac{802{,}5\times10^3}{337{,}1} \approx \SI{2680}{K}. \]
La valeur mesurée est plus proche de $\SI{2200}{K}$ : les capacités thermiques augmentent avec $T$, une partie de \ce{CO2} et de \ce{H2O} se dissocie (réactions endothermiques) et il y a des pertes. Le modèle donne néanmoins le bon ordre de grandeur. Le diazote, qui « ne fait rien », absorbe près de 70\,\% de l'énergie : c'est pourquoi une flamme dans \ce{O2} pur est beaucoup plus chaude (exercice~\ref{exo:gaz}).
\end{exemple}

\section{Capacité numérique : simulation d'un réacteur adiabatique}

Si la vitesse de la réaction est connue, on peut suivre $T(t)$. La vitesse augmente avec la température (loi d'Arrhenius $k(T) = A\,\mathrm{e}^{-E_a/RT}$) : une réaction exothermique en réacteur adiabatique \emph{s'auto-accélère}.

\begin{numerique}[modèle]
Réaction \ce{A -> B} d'ordre 1, dans un réacteur adiabatique monobare de capacité thermique totale $C$ :
\[ \frac{\dd\xi}{\dd t} = k(T)\,(n_0-\xi), \qquad C\,\dd T + \Delta_rH^\circ\,\dd\xi = 0 \;\Longrightarrow\; \frac{\dd T}{\dd t} = -\frac{\Delta_rH^\circ}{C}\,\frac{\dd\xi}{\dd t}. \]
Le bilan d'enthalpie entre $t$ et $t+\dd t$ est exactement celui du chemin fictif, appliqué à chaque pas de temps. On résout ce système par la méthode d'Euler explicite.
\end{numerique}

\begin{tcolorbox}[enhanced,breakable,colback=black!3,colframe=black!25,boxrule=0.4pt,arc=1mm,left=2mm,right=2mm]
\lstinputlisting{code/reacteur_adiabatique.py}
\end{tcolorbox}

\begin{figure}[H]
\centering
\begin{tikzpicture}
\begin{axis}[cours,width=0.78\linewidth,height=0.42\linewidth,xlabel={$t$ (min)},ylabel={$T$ ($\si{\celsius}$)},
  xmin=0,xmax=60,ymin=20,ymax=70,axis y line*=left,
  legend style={at={(0.03,0.97)},anchor=north west}]
  \addplot[c3] table[x=tmin,y=TC]{data/adiab.dat}; \addlegendentry{$T(t)$, adiabatique}
  \addplot[c3,dashed,line width=0.8pt] coordinates {(0,65) (60,65)};
  \node[c3,font=\small,anchor=south east] at (axis cs:59,65) {$T_I + \Delta T_{\text{ad}} = \SI{65}{\celsius}$};
\end{axis}
\begin{axis}[cours,width=0.78\linewidth,height=0.42\linewidth,xmin=0,xmax=60,ymin=0,ymax=1.25,
  axis x line=none,axis y line*=right,ylabel={$\xi$ (mol)},grid=none,
  legend style={at={(0.97,0.03)},anchor=south east}]
  \addplot[c1] table[x=tmin,y=xi]{data/adiab.dat}; \addlegendentry{$\xi(t)$, adiabatique}
  \addplot[c1,dashed] table[x=tmin,y=xiiso]{data/adiab.dat}; \addlegendentry{$\xi(t)$, isotherme $\SI{25}{\celsius}$}
\end{axis}
\end{tikzpicture}
\caption{Simulation ($n_0 = \SI{1}{mol}$, $\Delta_rH^\circ = \SI{-120}{kJ.mol^{-1}}$, $C = \SI{3,0}{kJ.K^{-1}}$, $E_a = \SI{80}{kJ.mol^{-1}}$, $k(\SI{298}{K}) = \SI{2,0e-4}{s^{-1}}$). Le temps de demi-réaction passe de $\SI{58}{min}$ (isotherme) à $\SI{22}{min}$ (adiabatique).}
\label{fig:adiab}
\end{figure}

\noindent\textbf{Lecture de la simulation.} La température finale vaut $T_I + \Delta T_{\text{ad}}$ avec $\Delta T_{\text{ad}} = -n_0\Delta_rH^\circ/C = \SI{40}{K}$ : on retrouve la formule du chemin fictif. Entre-temps, l'échauffement multiplie la constante de vitesse par
\[ \frac{k(T_F)}{k(T_I)} = \exp\left[\frac{E_a}{R}\left(\frac{1}{T_I}-\frac{1}{T_F}\right)\right] \approx 45, \]
d'où l'allure en S de la courbe $T(t)$, caractéristique de l'\emph{emballement thermique}. C'est le phénomène que l'on cherche à éviter en génie chimique (refroidissement, dilution, ajout progressif des réactifs).
